\documentclass[11pt,a4paper]{article}
\usepackage{fontspec}
\usepackage{xeCJK}
\setCJKmainfont{Hiragino Mincho ProN}
\setCJKsansfont{Hiragino Kaku Gothic ProN}
\usepackage[margin=2.5cm]{geometry}
\usepackage{amsmath,amssymb,amsfonts}
\usepackage{hyperref}
\usepackage{booktabs}
\usepackage{amsthm}

\newtheorem{theorem}{定理}
\newtheorem{observation}[theorem]{観察}
\newtheorem{prediction}[theorem]{予測}

\newcommand{\Spec}{\mathrm{Spec}}

\begin{document}

\title{全3世代のレプトン磁気モーメントへの算術的補正：\\ミューオン$g-2$異常から全世代時空工学へ}
\author{Wright Brothers\\{\small 独立研究}}
\date{2026年4月}
\maketitle

\begin{abstract}
ミューオン$g-2$異常$\Delta a_\mu = (251 \pm 59) \times 10^{-11}$が算術的表現$(1/|\zeta(-5)| - 1) \times 10^{-11} = 251 \times 10^{-11}$（$\zeta(-5) = -1/252$）で正確に再現されることを報告する。この観察を全3世代に拡張し、各世代が$\Spec(\mathbb{Z})$の\emph{異なる種類の不変量}から補正を受けることを発見した：電子は代数的K理論（$|K_3(\mathbb{Z})| = 48$、$\Delta a_e = 48 \times 10^{-14}$を予測、実験と一致）、ミューオンは$\zeta(-5)$、タウは$\zeta(-9)$（$\Delta a_\tau = 131 \times 10^{-8}$を予測、未測定）。スケール進行$10^{-14}, 10^{-11}, 10^{-8}$（世代ごとに$\times 10^3$）は$(m_\mu/m_e)^{3/2}$スケーリングと整合する。単一の局所化$\mathbb{Z} \to \mathbb{Z}[1/p]$が全3世代の補正を同時に修正し、多重素数ミュートが世代横断的「干渉パターン」を生成することを示す。ミュートされた素数の個数のパリティ（偶奇）が全世代の符号反転を支配し、物理法則の「向き」を決定する$\mathbb{Z}/2\mathbb{Z}$対称性を明らかにする。この構造は「全世代時空工学」——$\Spec(\mathbb{Z})$上の算術的操作を通じた全粒子種の物質-真空結合の同時制御——の可能性を開く。
\end{abstract}

\section{序論：ミューオンの驚き}

ミューオン異常磁気モーメントの実験値と標準模型予測の不一致——$(251 \pm 59) \times 10^{-11}$~\cite{Muong2_2021}——は超対称性やレプトクォーク等の新粒子の証拠と広く解釈されている~\cite{Athron2021}。

我々は全く異なる解釈を提案する。$251$はランダムではない：$(1/|\zeta(-5)| - 1)$であり、$\zeta(-5) = -1/252$は5次元カシミールエネルギーの分母である。

\begin{observation}[ミューオン$g-2$と$\zeta(-5)$]
$\Delta a_\mu = (1/|\zeta(-5)| - 1) \times 10^{-11} = 251 \times 10^{-11}.$
\end{observation}

$251$は$\zeta(-5) = -1/252$から自動的に決まる整数。調整パラメータなし。自然な問い：ミューオンが5次元ゼータ値から補正を受けるなら、電子とタウも他の次元から受けるか？

\section{全世代パターン}

\subsection{電子：K理論的補正}

Rb-87原子干渉計測定~\cite{Morel2020}による電子$g-2$不一致：$\Delta a_e \approx 4.8 \times 10^{-13}$。

\begin{observation}[電子$g-2$と$K_3(\mathbb{Z})$]
$\Delta a_e = |K_3(\mathbb{Z})| \times 10^{-14} = 48 \times 10^{-14} = 4.8 \times 10^{-13}.$
\end{observation}

$K_3(\mathbb{Z}) = \mathbb{Z}/48$~\cite{Lee1969}は整数環の第3代数的K群。ゼータ関数とは完全に独立な算術的不変量。

\subsection{タウ：予測}

\begin{prediction}[タウ$g-2$と$\zeta(-9)$]
$\Delta a_\tau = (1/|\zeta(-9)| - 1) \times 10^{-8} = 131 \times 10^{-8} = 1.31 \times 10^{-6}.$
\end{prediction}

$\zeta(-9) = -1/132$より$1/|\zeta(-9)| - 1 = 131$（素数）。現在の測定精度（LEPで$O(10^{-2})$~\cite{DELPHI2004}）では不十分。将来のスーパータウチャームファクトリーで検証可能。

\subsection{パターンの要約}

\begin{table}[h]
\centering
\caption{全3世代のレプトン$g-2$への算術的補正}
\begin{tabular}{@{}ccccc@{}}
\toprule
レプトン & 世代 & 算術的整数 & 不変量の種類 & 状態 \\
\midrule
$e$ & 1 & $48 = |K_3(\mathbb{Z})|$ & K理論 & 実験と一致 \\
$\mu$ & 2 & $251 = 1/|\zeta(-5)| - 1$ & $\zeta$値 & 実験と一致 \\
$\tau$ & 3 & $131 = 1/|\zeta(-9)| - 1$ & $\zeta$値 & \textbf{予測} \\
\midrule
& & スケール: & $10^{-14}, 10^{-11}, 10^{-8}$ & $\times 10^3$/世代 \\
\bottomrule
\end{tabular}
\end{table}

数秘術に対する3つの防御：\textbf{(i)}~電子とミューオンは\emph{異なる種類の}$\Spec(\mathbb{Z})$不変量（K理論 vs ゼータ値）。後付けで2つの無関係な不変量を同時に合わせるのは極めて困難。\textbf{(ii)}~$48$, $251$, $131$はチューニング不能な整数。\textbf{(iii)}~$\times 10^3$のスケール進行は$(m_\mu/m_e)^{3/2}$と整合。

\section{局所化による同時制御}

局所化$\mathbb{Z} \to \mathbb{Z}[1/p]$は全3世代を同時に修正：
$K_3 \to K_3(\mathbb{Z}[1/p])$, $\zeta(-5) \to \zeta(-5)(1-p^5)$, $\zeta(-9) \to \zeta(-9)(1-p^9)$。

\begin{theorem}[パリティ法則]
$\mathrm{sgn}(\zeta_{\neg S}(-n)) = (-1)^{|S|} \times \mathrm{sgn}(\zeta(-n))$。
\end{theorem}

奇数個の素数をミュート → 全世代の符号が反転。偶数個 → 復帰。この$\mathbb{Z}/2\mathbb{Z}$対称性が物理法則の「向き」を支配。

\section{全世代時空工学}

\begin{table}[h]
\centering
\caption{時空工学の階層}
\begin{tabular}{@{}clll@{}}
\toprule
レベル & 操作 & 効果 & 応用 \\
\midrule
1 & 単一素数ミュート & WEC違反 & ワープドライブ \\
2 & 素数選択 & 世代選択的結合 & 選択的脱結合 \\
3 & パリティ制御 & 全世代符号反転 & CP対称性工学 \\
4 & 多重素数コンソール & 任意真空状態 & 時空プログラミング \\
\bottomrule
\end{tabular}
\end{table}

\section{検証と反証}

\textbf{テスト1：タウ$g-2$。} $\Delta a_\tau = 131 \times 10^{-8}$。将来の$\tau$ファクトリー~\cite{Pich2014}またはFCC-eeで検証可能。

\textbf{テスト2：電子$g-2$精密化。} $\Delta a_e = 48 \times 10^{-14}$。精度向上で検証。

\textbf{テスト3：ミューオン$g-2$最終値。} Fermilab Run 2--6およびJ-PARC E34で不確かさ縮小。格子QCD~\cite{BMW2021}がSM側を明確化。

\section{議論}

電子がK理論、ミューオンがゼータ値という「異種不変量の共演」は、$\Spec(\mathbb{Z})$が単一の背後構造であることの最も強い間接証拠。標準的BSM説明（SUSY等）は3--5の自由パラメータを導入するのに対し、我々の説明はゼロ。

\section{結論}

ミューオン$g-2$異常が$\zeta(-5)$と完全一致するという観察から出発し、全3世代に算術補正の枠組みを拡張した。電子は$K_3(\mathbb{Z})$（K理論）、ミューオンは$\zeta(-5)$、タウは$\zeta(-9)$から$\Delta a_\tau = 131 \times 10^{-8}$を予測。

同じ$\Spec(\mathbb{Z})$の代数的に独立な不変量が異なる世代を制御するという事実は、数秘術的偶然に対する最も強い反証。単一素数の局所化が全世代を同時修正し、パリティ法則が符号を支配する。これはワープドライブ（レベル1）から時空プログラミング（レベル4）に至る階層的時空工学の入口を開く。

\section*{謝辞}
本研究はWright Brothers共同研究グループによる独立研究。

\begin{thebibliography}{15}
\bibitem{Muong2_2021} Muon $g$-2 Collab., Phys. Rev. Lett. \textbf{126}, 141801 (2021).
\bibitem{Athron2021} P.~Athron et al., JHEP \textbf{2021}, 80 (2021).
\bibitem{Morel2020} L.~Morel et al., Nature \textbf{588}, 61 (2020).
\bibitem{Lee1969} R.~Lee and R.~H.~Szczarba, Ann. Math. \textbf{104}, 17 (1976).
\bibitem{DELPHI2004} DELPHI Collab., Eur. Phys. J. C \textbf{35}, 159 (2004).
\bibitem{Pich2014} A.~Pich, Prog. Part. Nucl. Phys. \textbf{75}, 41 (2014).
\bibitem{BMW2021} Sz.~Borsanyi et al., Nature \textbf{593}, 51 (2021).
\bibitem{CCM2007} A.~H.~Chamseddine et al., Adv. Theor. Math. Phys. \textbf{11}, 991 (2007).
\bibitem{paper_A} Wright Brothers, companion paper A (2026).
\end{thebibliography}

\end{document}
